\section{Finite orders}

\import{set/finite.tex}
\import{order/bounds.tex}
\import{order/chain.tex}

\begin{proposition}\label{greatest_element_cons}
    Suppose $R$ is transitive and $R$ is connex on $A$.
    Suppose $Y\subseteq A$ and $m$ is a greatest element of $Y$ with respect to $R$.
    Suppose $x\in A$.
    Then there exists $p\in\cons{x}{Y}$ such that
        $p$ is a greatest element of $\cons{x}{Y}$ with respect to $R$.
\end{proposition}
\begin{proof}
    Follows by \auto{greatest_element,connex,transitive,cons_iff,elem_subseteq}.
\end{proof}

\begin{proposition}\label{least_element_cons}
    Suppose $R$ is transitive and $R$ is connex on $A$.
    Suppose $Y\subseteq A$ and $m$ is a least element of $Y$ with respect to $R$.
    Suppose $x\in A$.
    Then there exists $p\in\cons{x}{Y}$ such that
        $p$ is a least element of $\cons{x}{Y}$ with respect to $R$.
\end{proposition}
\begin{proof}
    Follows by \auto{least_element,connex,transitive,cons_iff,elem_subseteq}.
\end{proof}

\begin{proposition}\label{minimal_element_cons}
    Suppose $R$ is an order on $A$.
    Suppose $Y\subseteq A$ and
        $m$ is a minimal element of $Y$ with respect to $R$.
    Suppose $x\in A$.
    Then there exists $p\in\cons{x}{Y}$ such that
        $p$ is a minimal element of $\cons{x}{Y}$ with respect to $R$.
\end{proposition}
\begin{proof}
    $R$ is transitive and $R$ is antisymmetric by definition.
    \begin{byCase}
    \caseOf{$x\mathrel{R}m$.}
        $x\in\cons{x}{Y}$ by \obv{cons_left}.
        Show for all $y\in\cons{x}{Y}$ such that $y\mathrel{R}x$
            we have $y=x$.
        \begin{subproof}
            Fix $y\in\cons{x}{Y}$.
            Assume $y\mathrel{R}x$.
            $y=x$ or $y\in Y$ by \auto{cons_iff}.
            \begin{byCase}
            \caseOf{$y=x$.}
                Follows.
            \caseOf{$y\in Y$.}
                $y\mathrel{R}m$ by \obv{transitive}.
                Hence $y=m$ by \obv{minimal_element}.
                Thus $m\mathrel{R}x$.
                Hence $x=m$ by \obv{antisymmetric}.
            \end{byCase}
        \end{subproof}
        Hence $x$ is a minimal element of $\cons{x}{Y}$
            with respect to $R$ by \obv{minimal_element}.
    \caseOf{$x\not\mathrel{R}m$.}
        $m\in Y$ by \obv{minimal_element}.
        Hence $m\in\cons{x}{Y}$ by \auto{cons_iff}.
        Show for all $y\in\cons{x}{Y}$ such that $y\mathrel{R}m$
            we have $y=m$.
        \begin{subproof}
            Fix $y\in\cons{x}{Y}$.
            Assume $y\mathrel{R}m$.
            $y=x$ or $y\in Y$ by \auto{cons_iff}.
            \begin{byCase}
            \caseOf{$y=x$.}
                Contradiction.
            \caseOf{$y\in Y$.}
                Follows by \obv{minimal_element}.
            \end{byCase}
        \end{subproof}
        Hence $m$ is a minimal element of $\cons{x}{Y}$
            with respect to $R$ by \obv{minimal_element}.
    \end{byCase}
\end{proof}

\begin{proposition}\label{maximal_element_cons}
    Suppose $R$ is an order on $A$.
    Suppose $Y\subseteq A$ and
        $m$ is a maximal element of $Y$ with respect to $R$.
    Suppose $x\in A$.
    Then there exists $p\in\cons{x}{Y}$ such that
        $p$ is a maximal element of $\cons{x}{Y}$ with respect to $R$.
\end{proposition}
\begin{proof}
    $R$ is transitive and $R$ is antisymmetric by definition.
    \begin{byCase}
    \caseOf{$m\mathrel{R}x$.}
        $x\in\cons{x}{Y}$ by \obv{cons_left}.
        Show for all $y\in\cons{x}{Y}$ such that $x\mathrel{R}y$
            we have $y=x$.
        \begin{subproof}
            Fix $y\in\cons{x}{Y}$.
            Assume $x\mathrel{R}y$.
            $y=x$ or $y\in Y$ by \auto{cons_iff}.
            \begin{byCase}
            \caseOf{$y=x$.}
                Follows.
            \caseOf{$y\in Y$.}
                $m\mathrel{R}y$ by \obv{transitive}.
                Hence $y=m$ by \obv{maximal_element}.
                Thus $x\mathrel{R}m$.
                Hence $x=m$ by \obv{antisymmetric}.
            \end{byCase}
        \end{subproof}
        Hence $x$ is a maximal element of $\cons{x}{Y}$
            with respect to $R$ by \obv{maximal_element}.
    \caseOf{$m\not\mathrel{R}x$.}
        $m\in Y$ by \obv{maximal_element}.
        Hence $m\in\cons{x}{Y}$ by \auto{cons_iff}.
        Show for all $y\in\cons{x}{Y}$ such that $m\mathrel{R}y$
            we have $y=m$.
        \begin{subproof}
            Fix $y\in\cons{x}{Y}$.
            Assume $m\mathrel{R}y$.
            $y=x$ or $y\in Y$ by \auto{cons_iff}.
            \begin{byCase}
            \caseOf{$y=x$.}
                Contradiction.
            \caseOf{$y\in Y$.}
                Follows by \obv{maximal_element}.
            \end{byCase}
        \end{subproof}
        Hence $m$ is a maximal element of $\cons{x}{Y}$
            with respect to $R$ by \obv{maximal_element}.
    \end{byCase}
\end{proof}

\begin{theorem}\label{finite_greatest_element}
    Suppose $R$ is transitive and $R$ is connex on $A$.
    Suppose $X\subseteq A$ and $X$ is a finite set and $X$ is inhabited.
    Then there exists $m\in X$ such that
        $m$ is a greatest element of $X$ with respect to $R$.
\end{theorem}
\begin{proof}
    Let $C=\{Y\in\pow{X}\mid
        \text{$Y=\emptyset$ or there exists $m\in Y$ such that
            $m$ is a greatest element of $Y$ with respect to $R$}\}$.
    $\emptyset\in C$ by \obv{powerset_bottom}.

    Show that for all $Y,y$ such that $Y\in C$ and $y\in X$ we have
        $\cons{y}{Y}\in C$.
    \begin{subproof}
        Fix $Y$.
        Fix $y$.
        Assume $Y\in C$ and $y\in X$.
        $Y\subseteq X$ by \obv{pow_iff,subseteq}.
        Thus $Y\subseteq A$ by \obv{subseteq_transitive}.
        \begin{align*}
            y
                &\in X
                \\
                &\subseteq A
        \end{align*}
        $\cons{y}{Y}\subseteq X$ by \obv{cons_subseteq_intro}.
        Thus $\cons{y}{Y}\in\pow{X}$ by \obv{pow_iff,subseteq}.
        \begin{byCase}
        \caseOf{$Y=\emptyset$.}
            $y$ is a greatest element of $\cons{y}{Y}$ with respect to $R$
                by \auto{greatest_element,cons_iff,emptyset}.
            Thus there exists $m\in\cons{y}{Y}$ such that
                $m$ is a greatest element of $\cons{y}{Y}$ with respect to $R$
                by \obv{cons_left}.
            Follows by \auto.
        \caseOf{$Y\neq\emptyset$.}
            Take $m\in Y$ such that
                $m$ is a greatest element of $Y$ with respect to $R$.
            Thus there exists $p\in\cons{y}{Y}$ such that
                $p$ is a greatest element of $\cons{y}{Y}$ with respect to $R$
                by \obv{greatest_element_cons}.
            Follows by \auto.
        \end{byCase}
    \end{subproof}
    $X\in C$ by \obv{finite_induction}.
    $X\neq\emptyset$ by \obv{emptyset}.
\end{proof}

\begin{theorem}\label{finite_least_element}
    Suppose $R$ is transitive and $R$ is connex on $A$.
    Suppose $X\subseteq A$ and $X$ is a finite set and $X$ is inhabited.
    Then there exists $m\in X$ such that
        $m$ is a least element of $X$ with respect to $R$.
\end{theorem}
\begin{proof}
    Let $C=\{Y\in\pow{X}\mid
        \text{$Y=\emptyset$ or there exists $m\in Y$ such that
            $m$ is a least element of $Y$ with respect to $R$}\}$.
    $\emptyset\in C$ by \obv{powerset_bottom}.

    Show that for all $Y,y$ such that $Y\in C$ and $y\in X$ we have
        $\cons{y}{Y}\in C$.
    \begin{subproof}
        Fix $Y$.
        Fix $y$.
        Assume $Y\in C$ and $y\in X$.
        $Y\subseteq X$ by \obv{pow_iff,subseteq}.
        Thus $Y\subseteq A$ by \obv{subseteq_transitive}.
        $y\in A$ by \obv{elem_subseteq}.
        $\cons{y}{Y}\subseteq X$ by \obv{cons_subseteq_intro}.
        Thus $\cons{y}{Y}\in\pow{X}$ by \obv{pow_iff,subseteq}.
        \begin{byCase}
        \caseOf{$Y=\emptyset$.}
            $y$ is a least element of $\cons{y}{Y}$ with respect to $R$
                by \auto{least_element,cons_iff,emptyset}.
            Thus there exists $m\in\cons{y}{Y}$ such that
                $m$ is a least element of $\cons{y}{Y}$ with respect to $R$
                by \obv{cons_left}.
            Follows by \auto.
        \caseOf{$Y\neq\emptyset$.}
            Take $m\in Y$ such that
                $m$ is a least element of $Y$ with respect to $R$.
            Thus there exists $p\in\cons{y}{Y}$ such that
                $p$ is a least element of $\cons{y}{Y}$ with respect to $R$
                by \obv{least_element_cons}.
            Follows by \auto.
        \end{byCase}
    \end{subproof}
    $X\in C$ by \obv{finite_induction}.
    $X\neq\emptyset$ by \obv{emptyset}.
\end{proof}

\begin{theorem}\label{finite_minimal_element}
    Suppose $R$ is an order on $A$.
    Suppose $X\subseteq A$ and $X$ is a finite set and $X$ is inhabited.
    Then there exists $m\in X$ such that
        $m$ is a minimal element of $X$ with respect to $R$.
\end{theorem}
\begin{proof}
    Let $C=\{Y\in\pow{X}\mid
        \text{$Y=\emptyset$ or there exists $m\in Y$ such that
            $m$ is a minimal element of $Y$ with respect to $R$}\}$.
    $\emptyset\in C$ by \obv{powerset_bottom}.

    Show that for all $Y,y$ such that $Y\in C$ and $y\in X$ we have
        $\cons{y}{Y}\in C$.
    \begin{subproof}
        Fix $Y$.
        Fix $y$.
        Assume $Y\in C$ and $y\in X$.
        $Y\subseteq X$ by \obv{pow_iff,subseteq}.
        Thus $Y\subseteq A$ by \obv{subseteq_transitive}.
        $y\in A$ by \obv{elem_subseteq}.
        $\cons{y}{Y}\subseteq X$ by \obv{cons_subseteq_intro}.
        Thus $\cons{y}{Y}\in\pow{X}$ by \obv{pow_iff,subseteq}.
        \begin{byCase}
        \caseOf{$Y=\emptyset$.}
            $y$ is a minimal element of $\cons{y}{Y}$ with respect to $R$
                by \auto{minimal_element,cons_iff,emptyset}.
            Thus there exists $m\in\cons{y}{Y}$ such that
                $m$ is a minimal element of $\cons{y}{Y}$ with respect to $R$
                by \obv{cons_left}.
            Follows by \auto.
        \caseOf{$Y\neq\emptyset$.}
            Take $m\in Y$ such that
                $m$ is a minimal element of $Y$ with respect to $R$.
            Thus there exists $p\in\cons{y}{Y}$ such that
                $p$ is a minimal element of $\cons{y}{Y}$ with respect to $R$
                by \obv{minimal_element_cons}.
            Follows by \auto.
        \end{byCase}
    \end{subproof}
    $X\in C$ by \obv{finite_induction}.
    $X\neq\emptyset$ by \obv{emptyset}.
\end{proof}

\begin{theorem}\label{finite_maximal_element}
    Suppose $R$ is an order on $A$.
    Suppose $X\subseteq A$ and $X$ is a finite set and $X$ is inhabited.
    Then there exists $m\in X$ such that
        $m$ is a maximal element of $X$ with respect to $R$.
\end{theorem}
\begin{proof}
    Let $C=\{Y\in\pow{X}\mid
        \text{$Y=\emptyset$ or there exists $m\in Y$ such that
            $m$ is a maximal element of $Y$ with respect to $R$}\}$.
    $\emptyset\in C$ by \obv{powerset_bottom}.

    Show that for all $Y,y$ such that $Y\in C$ and $y\in X$ we have
        $\cons{y}{Y}\in C$.
    \begin{subproof}
        Fix $Y$.
        Fix $y$.
        Assume $Y\in C$ and $y\in X$.
        $Y\subseteq X$ by \obv{pow_iff,subseteq}.
        Thus $Y\subseteq A$ by \obv{subseteq_transitive}.
        $y\in A$ by \obv{elem_subseteq}.
        $\cons{y}{Y}\subseteq X$ by \obv{cons_subseteq_intro}.
        Thus $\cons{y}{Y}\in\pow{X}$ by \obv{pow_iff,subseteq}.
        \begin{byCase}
        \caseOf{$Y=\emptyset$.}
            $y$ is a maximal element of $\cons{y}{Y}$ with respect to $R$
                by \auto{maximal_element,cons_iff,emptyset}.
            Thus there exists $m\in\cons{y}{Y}$ such that
                $m$ is a maximal element of $\cons{y}{Y}$ with respect to $R$
                by \obv{cons_left}.
            Follows by \auto.
        \caseOf{$Y\neq\emptyset$.}
            Take $m\in Y$ such that
                $m$ is a maximal element of $Y$ with respect to $R$.
            Thus there exists $p\in\cons{y}{Y}$ such that
                $p$ is a maximal element of $\cons{y}{Y}$ with respect to $R$
                by \obv{maximal_element_cons}.
            Follows by \auto.
        \end{byCase}
    \end{subproof}
    $X\in C$ by \obv{finite_induction}.
    $X\neq\emptyset$ by \obv{emptyset}.
\end{proof}

\begin{theorem}\label{finite_subset_of_chain_union}
    Suppose $M$ is inhabited and $M$ is an inclusion chain.
    Suppose $F$ is a finite set and $F\subseteq\unions{M}$.
    Then there exists $E\in M$ such that $F\subseteq E$.
\end{theorem}
\begin{proof}
    Let $K=\{Y\in\pow{F}\mid
        \text{there exists $E\in M$ such that $Y\subseteq E$}\}$.
    Take $E_0$ such that $E_0\in M$.
    $\emptyset\subseteq E_0$ by \obv{emptyset_subseteq}.
    Thus $\emptyset\in K$ by \obv{powerset_bottom}.

    Show that for all $Y,y$ such that $Y\in K$ and $y\in F$ we have
        $\cons{y}{Y}\in K$.
    \begin{subproof}
        Fix $Y$.
        Fix $y$.
        Assume $Y\in K$ and $y\in F$.
        Take $C\in M$ such that $Y\subseteq C$.
        $y\in\unions{M}$ by \obv{elem_subseteq}.
        Take $D\in M$ such that $y\in D$ by \obv{unions_iff}.
        $C\subseteq D$ or $D\subseteq C$ by \obv{inclusion_chain}.
        $Y\subseteq F$ by \obv{pow_iff,subseteq}.
        Thus $\cons{y}{Y}\subseteq F$ by \obv{cons_subseteq_intro}.
        Hence $\cons{y}{Y}\in\pow{F}$ by \obv{pow_iff,subseteq}.
        \begin{byCase}
        \caseOf{$C\subseteq D$.}
            $Y\subseteq D$ by \obv{subseteq_transitive}.
            Thus $\cons{y}{Y}\subseteq D$ by \obv{cons_subseteq_intro}.
            Hence there exists $E\in M$ such that $\cons{y}{Y}\subseteq E$.
            Follows by \auto.
        \caseOf{$D\subseteq C$.}
            $y\in C$ by \obv{elem_subseteq}.
            Thus $\cons{y}{Y}\subseteq C$ by \obv{cons_subseteq_intro}.
            Hence there exists $E\in M$ such that $\cons{y}{Y}\subseteq E$.
            Follows by \auto.
        \end{byCase}
    \end{subproof}
    $F\in K$ by \obv{finite_induction}.
\end{proof}
